\documentclass{article}
\usepackage[T1]{fontenc}
\usepackage{lmodern}
\usepackage{graphicx}
\usepackage{amsmath,amssymb}
\usepackage[a4paper,margin=2cm]{geometry}
\usepackage[absolute,overlay]{textpos}
\setlength{\TPHorizModule}{1mm}
\setlength{\TPVertModule}{1mm}
\usepackage[indonesian]{babel}
\usepackage{hyperref}
\setlength{\emergencystretch}{2em}
% unit: Halaman 57

\begin{document}

% === Halaman 57 (python/native) ===

• Permutasi adalah teknik untuk mengacak susunan bit di dalam 

sebuah blok bit menghasilkan susunan baru.

• Pada gambar di atas, posisi bit-bit di dalam blok input 32-bit 

dipermutasi menjadi susunan baru. 

16

1

2

3

4

5

6

7

8

9

10

11

12

13

14

15

17

18

19

20

21

22

23

24

25

26

27

28

29

30

31

32

16

1

2

3

4

5

6

7

8

9

10

11

12

13

14

15

17

18

19

20

21

22

23

24

25

26

27

28

29

30

31

32

57

\end{document}
